\documentclass[../main.tex]{subfiles}
\graphicspath{{\subfix{../image/}}}
\begin{document}
\chapter{向量场}\label{chapter:ch08-向量场}

向量场是多变量微积分中熟悉的研究对象。在那里，开子集 $U\subseteq\mathbb{R}^n$ 上的向量场就是从 $U$ 到 $\mathbb{R}^n$ 的连续映射，可以想象为在 $U$ 的每一点附上一个"箭头"。本章我们说明如何将这一想法推广到光滑流形上。

我们将抽象光滑流形 $M$ 上的向量场视为从 $M$ 到其切丛的一类特殊连续映射 $X$——它为每一点 $p\in M$ 指定一个切向量 $X_p\in T_pM$。引入定义后，我们探讨向量场在光滑映射微分下的行为方式。

下一节引入 Lie 括号运算，这是将两个光滑向量场组合得到另一个向量场的方式。然后描述 Lie 括号最重要的应用：Lie 群上所有在左平移下不变的光滑向量场的集合在 Lie 括号下封闭，因而构成一个自然与该群相关联的代数对象，称为该 Lie 群的 Lie 代数。我们描述 Lie 代数的一些基本性质，并计算若干熟悉群的 Lie 代数。然后说明 Lie 群同态如何诱导其 Lie 代数的同态，由此可知同构的 Lie 群有同构的 Lie 代数。最后，在本章末尾说明如何识别 Lie 子群的 Lie 代数。

\subfile{Chapter-08/section-01.tex}
\subfile{Chapter-08/section-02.tex}
\subfile{Chapter-08/section-03.tex}
\subfile{Chapter-08/section-04.tex}
\subfile{Chapter-08/section-05.tex}
\end{document}